\section{源码等价判定关系}

本章只记录 \srcpath{src/sppt/metamorphic.cpp} 与
\srcpath{src/sppt/guard.cpp} 当前执行的判定。头文件
\srcpath{include/hacdcpf/sppt/metamorphic.hpp:MetamorphicObservation} 的注释列出 MR1--MR8，
但当前公共声明和实现中不存在 MR8 函数；因此本章不把“DAE 指标 1”写成已实现关系。

\subsection{公共差值算子}

\srcpath{max\_abs\_diff} 只比较两个向量公共前缀：
\begin{equation}
 d_{\max}(a,b)=\max_{0\le i<\min(|a|,|b|)}|a_i-b_i|,
\end{equation}
空公共前缀返回零；该函数自身不报告长度不一致，见
\srcpath{src/sppt/metamorphic.cpp:max\_abs\_diff}。
Ybus 差值先分别调用 \srcpath{make\_solver\_data\_projected}；维数不同返回无穷大，
零维返回零，其余返回稠密差矩阵元素绝对值最大值，见
\srcpath{src/sppt/metamorphic.cpp:ybus\_max\_abs\_diff}。

\subsection{MR1：投影幂等判定}

令 $p_1=\Pi(S)$、$p_2=\Pi(p_1)$。代码计算 AC/DC 母线数差绝对值之和
$d_b$、AC/DC 支路数差绝对值之和 $d_l$，以及负荷、发电机、换流器数量差
$d_{load},d_g,d_c$；另计算总有功负荷差 $d_P$、总无功负荷差 $d_Q$ 和
Ybus 最大差 $d_Y$。报告残差为
\begin{equation}
 r=d_b+d_l+d_{load}+d_g+d_c+d_P+d_Q+d_Y,
\end{equation}
但通过条件不是简单的 $r\le10^{-12}$，而是
\begin{equation}
 d_b=d_l=d_{load}=d_g=d_c=0,\quad d_P=d_Q=0,\quad d_Y\le10^{-12}.
\end{equation}
依据为 \srcpath{src/sppt/metamorphic.cpp:mr1\_projection\_idempotence}。

\subsection{MR2：归因往返判定}

无 \srcpath{bus\_merge\_map} 或映射为空时直接通过且残差为零。
否则在合并空间构造 $v_t=1+t$，按 Intensive 语义反投影；遍历
\fld{ext\_to\_int} 时跳过 dead bus、缺失原位置、越界原位置和越界内部位置，残差为
\begin{equation}
 r=\max_{e\in\mathcal E_{checked}}|v^{ext}_{pos(e)}-v^{int}_{map(e)}|,
\end{equation}
并以 $r\le\fld{tol}$ 通过，见 \srcpath{src/sppt/metamorphic.cpp:mr2\_attribution\_roundtrip}。

\subsection{MR3、MR3d 与 MR3t}

MR3 分别求解原系统和显式投影系统；任一未收敛即失败且残差固定为 1。
若两个电压向量分别匹配各自 AC 母线数，代码按稳定母线 ID 比较投影系统中能在原系统找到的
$V_m$ 和 $V_a$；否则只对投影结果的 $V_m$ 做 Intensive 反投影，长度仍不一致则残差固定为 1。
DC 电压仅在两个向量均非空且长度相同时加入公共前缀最大差。最终以最大差
$r\le\fld{tol}$ 通过，见 \srcpath{src/sppt/metamorphic.cpp:mr3\_semantic\_preservation\_pf}。

MR3d 同样先要求两个 DCOPF 均收敛。任一 LMP 向量为空时，代码设置
\fld{passed=false}、\fld{observation=NotObserved}、$r=0$；长度/母线数匹配时按 ID 比较，
否则长度相同则按位置比较，再否则尝试 Intensive 反投影，反投影后长度仍不同则失败且 $r=1$。
依据为 \srcpath{src/sppt/metamorphic.cpp:mr3\_semantic\_preservation\_opf\_dual}。

MR3t 固定设置 $t_{end}=0.10$、$\Delta t=0.005$，启用潮流初始化和一致动态初值。
任一仿真失败或缺少末快照时失败且 $r=1$。否则残差取末快照最大/最小 AC 电压、频率、
COI 频率四个绝对差，以及仅在尺寸相同且非空时加入的状态向量和 DC 电压向量最大差；
以 $r\le\fld{tol}$ 通过，见 \srcpath{src/sppt/metamorphic.cpp:mr3\_semantic\_preservation\_transient}。

\subsection{MR4--MR7}

MR4 在没有非平凡合并时直接通过。否则对每个成员数至少为 2 的组，逐成员读取
\fld{ext\_to\_int}，缺键视为 $-1$；与组内首成员内部索引不同即把违规数加一。
残差就是违规数，通过条件为违规数等于零；输入的 \fld{tol} 不参与通过判定，见
\srcpath{src/sppt/metamorphic.cpp:mr4\_merged\_bus\_equipotential}。

MR5 先求解原系统，再仅反转 AC 母线 vector 顺序后求解。任一未收敛或 $V_m$ 长度不等于相应
母线数时失败且 $r=1$；否则按母线稳定 ID 比较能匹配的电压幅值，以最大差不超过
\fld{tol} 通过，见 \srcpath{src/sppt/metamorphic.cpp:mr5\_relabel\_invariance}。

MR6 所谓“不相交并”只拼接 AC 母线、AC 支路、发电机、负荷、并联元件、DC 母线和 DC 支路
七类 vector，没有拼接其余组件。残差只统计投影后 AC/DC 母线数和支路数相对两个独立投影数量和
的绝对差；通过条件为两项均为零，见 \srcpath{src/sppt/metamorphic.cpp:mr6\_compositionality}。

MR7 将所有 AC SLACK 母线降为 PQ、清除发电机 \fld{is\_slack}、停运外部电网，并清除 VSC
构网角色；随后调用 SolverReady 校验。只要任一 Error 消息不区分大小写包含
\texttt{slack}、\texttt{reference} 或 \texttt{no swing} 即通过。其残差无论通过与否均为零，
见 \srcpath{src/sppt/metamorphic.cpp:mr7\_wellposedness\_gate}。

\subsection{核心套件与未实现关系}

\srcpath{run\_core\_metamorphic\_suite} 当前依次运行 MR1、MR2、MR3、MR3d、MR4、MR5、MR7；
它不运行需要两个系统输入的 MR6，也不存在 MR8 调用，见
\srcpath{src/sppt/metamorphic.cpp:mr7\_wellposedness\_gate}。因此“核心套件通过”不能解释为 MR1--MR8 全部通过。

\subsection{守卫判定}

默认请求 AC 母线电压；存在 DC 母线、VSC、开关或断路器时分别追加对应观测量。
守卫先调用 \srcpath{validate\_full}；异常时立即返回 validation 失败。
适定性条件为校验问题中不存在消息含参考关键词的 Error，且
\srcpath{evaluate\_converter\_coordination(sys,true)} 没有 blocking issue。
归因完整性要求每个请求观测量的 \srcpath{coverage.total()} 均为真；投影或归因抛异常时为假。
最终代码布尔式为
\begin{equation}
 \fld{accepted}=\fld{validation\_ok}\land\fld{well\_posed}\land\fld{attribution\_total}.
\end{equation}
依据为 \srcpath{src/sppt/guard.cpp:guard\_system}。
